\chapter{Literature Review}

\section{Personalised Driver Risk Assessment with Adaptive Feedback using Crowdsensed Telemetric Data}

The rapid evolution of Intelligent Transport Systems (ITS) has brought significant advancements in how driver behavior is monitored and evaluated. Traditional approaches often relied on expensive, proprietary in-vehicle hardware such as On-Board Diagnostics (OBD-II) dongles or dedicated inertial measurement units (IMUs). However, recent trends have shifted towards utilizing the ubiquitous presence of smartphones. This paper by Auwal Sagir Muhammad et al., published in IET Intelligent Transport Systems (IET, 2025), presents a comprehensive framework for assessing driver risk using such crowdsensed telemetric data. The study focuses on transforming raw GPS traces and temporal data into structured, actionable driving risk indicators. By avoiding complex hardware, this approach democratizes driver profiling, making it accessible to a wider population and enabling large-scale, crowdsourced data collection for traffic safety analysis.

\subsection{Methodology}
The authors utilized a vast dataset of crowdsensed telemetric data collected from standard smartphones to engineer several highly specific GPS-derived driver risk indicators. The core of their methodology involved developing algorithms to filter out GPS noise and smooth the trajectory data before analysis. From this refined data, they established specific numerical thresholds for critical driving events. For instance, harsh braking was calculated by analyzing negative acceleration thresholds over consecutive GPS points, while rapid acceleration was identified through sudden positive velocity spikes. Overspeeding was determined by continuously comparing the GPS-derived vehicle speed against mapped speed limits obtained from external routing APIs. Furthermore, the turning rate was computed by analyzing the delta in GPS bearing over time, identifying aggressive cornering. 

Crucially, the system incorporated contextual flags based on the driving environment. Nighttime driving and rainy weather conditions were used to dynamically modulate the risk thresholds, acknowledging that a safe speed in broad daylight might be inherently dangerous during a nighttime downpour. These engineered indicators provided a structured, rule-based approach for extracting interpretable driving risk features from raw telemetric data. Various machine learning models, including Support Vector Machines and Random Forests, were then applied to evaluate how these engineered indicators effectively capture driving volatility and predict overall driver risk profiles.

\subsection{Advantages}
The proposed system introduces validated feature engineering methods for transforming raw GPS, temporal, and weather information into highly meaningful driver behavior indicators. A significant advantage of this system is its inherent scalability and cost-effectiveness; by relying entirely on crowdsensed smartphone data, it eliminates the need for any specialized vehicle modifications. Furthermore, by relying on well-defined, threshold-based rules rather than opaque deep learning models for feature extraction, the approach provides an easily interpretable solution. This transparency allows the system to generate adaptive, understandable feedback for the driver, explaining exactly which behaviors (e.g., frequent harsh braking) contributed to their risk profile.

\subsection{Disadvantages}
Despite its strengths, the methodology presents notable limitations. One major limitation is that the proposed thresholds for harsh braking, acceleration, and turning were validated within a specific machine learning pipeline and a particular regional context. Driving cultures and road topologies vary drastically across the globe; therefore, these rigid numerical values may require significant recalibration when applied to different geographical driving environments. Additionally, a pure telematics approach lacks visual context. If the system detects harsh braking, it cannot determine whether the driver was driving aggressively or reacting appropriately to a child suddenly running into the street. This lack of visual situational awareness highlights the need for systems that integrate telematics with camera-based perception, a gap that the ContextDrive system aims to bridge.

\section{A New CNN Based Object Detection System for Autonomous Mobile Robots Based on Real World Vehicle Datasets}

The deployment of accurate and reliable object detection models is a cornerstone of modern autonomous navigation and advanced driver-assistance systems. While cloud-based deep learning models have achieved remarkable accuracy, they suffer from latency and require continuous internet connectivity, making them unsuitable for critical driving safety applications. This paper by Udink Aulia et al., published in Heliyon (Elsevier, 2024), addresses this challenge by presenting a lightweight object detection system specifically designed for autonomous mobile robots and edge devices. The study focuses on balancing real-time processing speeds with high detection accuracy in real-world environments, heavily utilizing transfer learning to adapt generic models to specific navigational tasks.

\subsection{Methodology}
To achieve real-time performance on hardware-constrained edge devices, the authors developed a vehicle detection system utilizing the SSD (Single Shot MultiBox Detector) framework paired with a MobileNetV2 backbone and an FPN-Lite (Feature Pyramid Network) architecture. The MobileNetV2 backbone employs depthwise separable convolutions to drastically reduce the computational complexity and parameter count compared to standard convolutional neural networks (CNNs). The FPN-Lite structure allows the model to construct high-level semantic feature maps at all scales, significantly improving the detection of small or distant vehicles, which is critical for early collision warning.

To overcome the limitations of generic pre-trained models, which often fail to generalize to specific deployment environments, the authors employed extensive transfer learning. The pre-trained detector was fine-tuned on a newly collected and meticulously annotated real-world vehicle dataset. This dataset was specifically curated to represent the lighting, weather, and traffic conditions the autonomous robots would actually face. Several training configurations, including various learning rates, batch sizes, and data augmentation techniques (such as random cropping, flipping, and color jittering), were tested. The models were then rigorously benchmarked to identify the optimal hyperparameters that maximize the mean Average Precision (mAP) while keeping inference latency low enough for real-time mobile robot navigation.

\subsection{Advantages}
The study provides compelling empirical evidence that transfer learning and fine-tuning of a lightweight SSD MobileNetV2 detector can significantly improve vehicle detection accuracy while maintaining strict real-time performance. This provides a highly practical, ready-to-deploy approach for domain-specific implementation on resource-constrained edge devices. By demonstrating that high accuracy does not necessitate massive computing power, the research paves the way for advanced perception systems to be deployed on affordable hardware, making safety technologies more accessible.

\subsection{Disadvantages}
The primary disadvantage of this approach is the phenomenon known as ``domain shift.'' The detection performance of the fine-tuned model heavily depends on the quality, size, and diversity of the fine-tuning dataset. If the autonomous robot or vehicle encounters conditions not well-represented in the training data---such as sudden heavy fog, snow, or nighttime driving in poorly lit areas---the model's accuracy degrades rapidly. Consequently, adapting the model to completely new environments requires substantial, time-consuming additional data collection and manual annotation efforts. This highlights the limitation of relying solely on a fixed visual dataset without integrating dynamic environmental context APIs to adjust system sensitivity.

\section{Using Contextual Data to Predict Risky Driving Events: A Novel Methodology from Explainable Artificial Intelligence}

Traditionally, driving risk assessment has focused heavily on isolated telematics variables, such as absolute speed or the frequency of hard braking. However, driving is an inherently contextual activity; a speed of 60 km/h might be perfectly safe on a dry highway but extremely dangerous on an icy residential street. This paper by Leandro Masello et al., published in Accident Analysis \& Prevention (Elsevier, 2023), investigates how these surrounding contextual factors fundamentally influence risky driving events. The research introduces a novel methodology based on explainable artificial intelligence (XAI) to move beyond isolated risk variables, analyzing how combined environmental and traffic contexts compound to create hazardous driving situations. By bridging the gap between raw data collection and human-understandable reasoning, this study highlights the growing need for transparency in modern intelligent transport systems.

\subsection{Methodology}
The researchers utilized large-scale, naturalistic telematics data collected from fleets of vehicles over an extended period. This comprehensive dataset included standard variables like GPS coordinates, vehicle speed, and acceleration profiles. Crucially, they augmented this telematics data with a rich array of contextual driving factors obtained from external geographical and meteorological databases. These included varying speed limits, real-time weather conditions (precipitation levels, visibility limits), traffic density metrics, road characteristics (such as urban vs. rural settings, intersection proximity, and road curvature), and temporal variables (time of day, day of week). 

They trained sophisticated supervised machine learning models, specifically tree-based ensembles like Random Forest and Gradient Boosting Machines (GBM), to predict the occurrence of risky driving events based on this multi-dimensional data. Tree-based models were chosen for their robust ability to handle non-linear relationships and heterogeneous data types without extensive preprocessing. However, to interpret the complex, non-linear interactions between these variables learned by the black-box models, they applied advanced XAI techniques, specifically SHAP (SHapley Additive exPlanations). Based on cooperative game theory, SHAP values allowed the researchers to quantify the exact marginal contribution of each individual contextual factor toward the likelihood of a risky driving event occurring. For instance, the SHAP analysis could reveal precisely how much "heavy rain" contributed to a model predicting a "high risk" event compared to "speeding." This provided a highly granular, feature-by-feature view of how risk is dynamically formed in real-time.

\subsection{Advantages}
The methodology effectively and quantitatively demonstrates that driving risk is determined by the combined, compounded influence of multiple contextual factors rather than isolated telematics variables. It provides a strong, evidence-based foundation for context-aware risk assessment. For example, the study proved that the danger of a sharp turn is exponentially increased when combined with poor visibility, a correlation that simple rule-based telematics systems often miss. 

Furthermore, the use of SHAP values transforms opaque machine learning predictions into transparent insights. This transparency is vital for the widespread adoption and psychological acceptance of driver assistance systems. By showing a driver exactly *why* a situation was classified as risky (e.g., ``High risk due to combination of moderate speed and heavy rain''), the system builds user trust and encourages active behavioral correction much more effectively than a generic, unexplained warning beep. It shifts the paradigm from a reactive warning system to a proactive, educational driving companion.

\subsection{Disadvantages}
Despite its significant contributions to explainability, the proposed methodology presents practical challenges for real-world deployment on edge devices. The system relies heavily on large-scale historical telematics data for supervised machine learning training. As a result, the models require extensive, computationally expensive retraining and recalibration when applied to completely different driving environments, varying vehicle types, or new geographical datasets where the baseline contextual risk factors might fundamentally differ.

Additionally, calculating SHAP values for complex ensemble models is incredibly resource-intensive. While this methodology is highly suitable for post-trip analysis, driver debriefing, or batch processing on powerful backend servers, running these XAI algorithms continuously in real-time on a resource-constrained edge device or consumer smartphone presents a massive computational bottleneck. For real-time applications like ContextDrive, simpler, deterministic rule-based engines or highly optimized, quantized neural networks must be employed to provide instant feedback without draining the device's battery or causing thermal throttling, even if they sacrifice some of the deep nuance provided by full SHAP analysis.
\section{Optimizing Monocular Driving Assistance for Real Time Processing on Jetson AGX Xavier}

While autonomous vehicles often rely on expensive, multi-modal sensor suites including LiDAR, radar, and stereo cameras, consumer-grade Advanced Driver Assistance Systems (ADAS) must balance performance with cost and portability. Monocular camera systems offer a highly affordable alternative but require sophisticated computer vision algorithms to extract depth and spatial awareness from a single lens. This paper by Huy-Hung Nguyen et al., published in IEEE Access (IEEE, 2024), presents a monocular camera-based ADAS optimized specifically for high-performance edge computing, aiming to perform complex road scene understanding in real time without specialized depth sensors.

\subsection{Methodology}
The authors developed a unified vision pipeline designed to simultaneously handle multiple perception tasks, such as vehicle detection, pedestrian tracking, and traffic sign recognition, using a single live monocular camera feed. The core of their methodology involved deeply optimizing these lightweight object detection models for deployment on the NVIDIA Jetson AGX Xavier, a powerful embedded edge computing platform. 

To achieve high frame rates, they integrated aggressive optimization techniques. This included tensor optimization using NVIDIA's TensorRT, which fuses network layers and optimizes memory access patterns. Furthermore, they employed reduced precision inference, converting the neural network weights from standard 32-bit floating-point (FP32) to 16-bit floating-point (FP16) or 8-bit integer (INT8) representations. This quantization drastically reduces memory bandwidth requirements and accelerates mathematical operations. The performance of these highly optimized models was then rigorously benchmarked against standard baseline architectures to meticulously evaluate the trade-offs between frame rate throughput, processing latency, and overall detection accuracy in highly dynamic driving scenarios.

\subsection{Advantages}
The research convincingly demonstrates that a single monocular camera, when paired with heavily optimized lightweight object detection algorithms, can reliably perform complex, real-time perception tasks required for driver assistance. By achieving high frame rates and low latency, the system validates the feasibility of comprehensive camera-based road scene understanding without needing specialized, expensive sensors like LiDAR. This approach significantly reduces the hardware cost and complexity of implementing ADAS features in standard vehicles.

\subsection{Disadvantages}
The primary limitation of this research lies in its hardware specificity. The system's architecture, memory management, and specific tensor optimization enhancements are heavily tailored specifically for the NVIDIA Jetson AGX Xavier hardware platform and its unique GPU architecture. While powerful, the Jetson AGX Xavier is a relatively expensive embedded system not typically found in consumer smartphones. Consequently, the deployment strategy and acceleration techniques developed in this study are not directly transferable to standard smartphone-based implementations, which rely on different neural processing units (NPUs) and frameworks like Google LiteRT (TensorFlow Lite), requiring entirely different optimization paradigms.

\section{Balancing Transparency and Control: The Impact of AI Explanation Detail on User Perception in Automated Vehicles}

As artificial intelligence takes a profoundly larger role in vehicle control and driver assistance, the human-computer interaction (HCI) aspect of these systems has emerged as a critical bottleneck for widespread public adoption. A significant, persistent barrier to the acceptance of AI-driven safety systems is the well-documented ``black box'' problem, where users do not understand the underlying mechanisms dictating why a system made a certain sudden decision. This opacity often leads to ``automation surprise,'' resulting in either a complete lack of trust and system disabling, or conversely, dangerous over-reliance. This paper by Jakob Peintner et al., published in Transportation Research Interdisciplinary Perspectives (Elsevier, 2025), rigorously investigates the psychological and behavioral effects of AI explainability on user trust and perception. The core objective of the study explores how varying levels of detail and complexity in safety explanations affect a user's cognitive experience, aiming to pinpoint the optimal balance between providing necessary transparency to build trust and preventing cognitive overload that could impair situational awareness.

\subsection{Methodology}
To safely, ethically, and consistently evaluate user reactions to critical, potentially dangerous driving events, the authors designed and conducted a comprehensive virtual reality (VR) driving study. Human participants were placed in a highly realistic, immersive simulated Level-4 highly automated vehicle environment. In this setup, the vehicle handles the vast majority of the dynamic driving tasks, placing the human participant primarily in a supervisory or passenger role. Throughout the controlled simulation, the vehicle deliberately encountered various complex, randomized traffic scenarios requiring sudden automated interventions, such as a pedestrian unexpectedly darting into the road or a lead vehicle braking abruptly.

To isolate the effect of explainability, the researchers systematically tested four distinct AI explanation levels presented via an in-cabin graphical interface during these critical interventions: ``None'' (the vehicle acts completely autonomously without providing any rationale), ``How'' (the system simply describes the physical action it is taking, e.g., ``Braking hard''), ``How \& Why'' (the system describes both the action and its underlying immediate reasoning, e.g., ``Braking hard because of a pedestrian ahead''), and finally ``Contrastive \& Counterfactual'' (the system exhaustively explains why an alternative action was inherently worse and not taken, e.g., ``Braking hard because swerving into the left lane would cause a catastrophic collision with oncoming traffic''). 

To capture a holistic view of the user experience, the researchers employed a multi-modal data collection approach. Beyond standard detailed post-scenario psychometric surveys (utilizing validated Trust in Automation scales), they continuously monitored physiological and behavioral metrics during the VR simulation. This included advanced eye-tracking to measure pupil dilation as a proxy for cognitive load and stress levels, allowing them to quantitatively evaluate the psychological toll of processing each explanation level during an emergency.

\subsection{Advantages}
The study provides robust, empirical evidence addressing a core HCI challenge in vehicle automation, shifting the discourse from anecdotal UI design to data-driven psychology. The research demonstrates clearly and quantitatively that concise ``How \& Why'' explanations offer the absolute most effective balance between system transparency and user usability. The physiological data proved that while providing no explanation severely degrades trust and spikes anxiety (due to automation surprise), providing overly complex, paragraph-style contrastive explanations drastically increases cognitive load without yielding any significant improvement in perceived safety or trust. 

This research establishes a clear, evidence-based psychological framework for designing understandable and trustworthy UI/UX for driver assistance recommendations. It validates the design philosophy that ADAS interfaces must immediately answer the user's instinctive ``Why is the car doing this?'' question concisely. This perfectly aligns with modern, context-aware systems like ContextDrive, which strive to distill complex risk calculations into immediate, digestible visual feedback (such as color-coded risk rings paired with short, specific ``Why'' and ``Advice'' strings) to build rapid user consensus.

\subsection{Disadvantages}
Despite the rigorous experimental design, the primary limitation of this study is its ecological validity regarding active driver assistance. The findings are strictly based on passenger interactions within a simulated Level-4 highly automated driving environment, where the user is essentially a passive observer with the luxury of time to process graphical notifications. 

Consequently, these results may not directly, safely generalize to real-time, manual driver assistance systems (Level 1/2) where the human user is actively controlling the vehicle's steering and pedals, and requires split-second, visceral alerts to avoid collisions. In highly dynamic, high-stakes manual driving situations, even reading a relatively concise textual ``How \& Why'' explanation on a smartphone screen might cause a critical visual distraction, drawing the driver's eyes away from the road precisely when their attention is needed most. For active real-time ADAS, explanations must be delivered through highly intuitive, non-intrusive modalities---such as specific, distinct audio tones, haptic feedback, or peripheral color changes---rather than relying primarily on explicit textual descriptions, a crucial real-world nuance that this passenger-centric VR study does not fully encompass.
